% !TEX root = ../../tree_notes.tex
% =====================================================================
%  Chapter 1 review -- summary, checklist, quick reference, mixed set,
%  related topics. Uses \chaptersection (unnumbered) and NO numbered boxes
%  (they would continue the 1.4.x numbering).
% =====================================================================

\chaptersection{Chapter 1 Summary}

Everything in this chapter is a tool the rest of the book assumes. The table collects
the results with the exact conditions they need.

\begin{center}
\renewcommand{\arraystretch}{1.25}
\begin{tabular}{@{}p{0.30\linewidth} p{0.46\linewidth} p{0.16\linewidth}@{}}
\toprule
\textbf{Result} & \textbf{Statement (with conditions)} & \textbf{Where} \\
\midrule
Edges of a tree    & A tree on $n$ nodes has exactly $n-1$ edges. & \thmref{c01-edges-nodes} \\
Unique path        & Exactly one simple path joins any two nodes (connected, acyclic). & \thmref{c01-unique-path} \\
Height recursion   & $\height(\nil)=-1$, $\height(v)=1+\max_c \height(c)$. & \thmref{c01-height-recursion} \\
Size recursion     & $\size(\nil)=0$, $\size(v)=1+\sum_c \size(c)$. & \thmref{c01-height-recursion} \\
Perfect counts     & Height $h$ (edges) $\Rightarrow$ $2^{h+1}-1$ nodes, $2^{h}$ leaves. & \thmref{c01-perfect-count} \\
Full tree leaves   & Every internal node has 2 children $\Rightarrow L = I+1$. & \thmref{c01-full-binary-leaves} \\
Height bounds      & $\floor{\log_2 n}\le \height \le n-1$ (binary, $n$ nodes). & \corref{c01-height-bounds} \\
Catalan shapes     & Binary shapes on $n$ nodes $= C_n = \frac{1}{n+1}\binom{2n}{n}$. & \thmref{c01-catalan-shapes} \\
Cayley trees       & Labelled (free) trees on $n$ vertices $= n^{\,n-2}$. & \thmref{c01-cayley} \\
\bottomrule
\end{tabular}
\end{center}

\medskip
{\small\textbf{Algorithms from this chapter.}}
\begin{center}
\renewcommand{\arraystretch}{1.2}
\begin{tabular}{@{}p{0.34\linewidth} l l l@{}}
\toprule
\textbf{Algorithm} & \textbf{Time} & \textbf{Memory} & \textbf{Code} \\
\midrule
Build adjacency from edges            & $O(n)$     & $O(n)$      & \coderef{c01-build} \\
One DFS: depth, height, size          & $O(n)$     & $O(n)$      & \coderef{c01-measure} \\
Structural equality (Same Tree)       & $O(\min(n,m))$ & $O(h)$ & \coderef{c01-same-tree} \\
Balanced check (single pass)          & $O(n)$     & $O(h)$      & \coderef{c01-balanced} \\
Count complete-tree nodes             & $O(\log^2 n)$ & $O(\log n)$ & \coderef{c01-count-nodes} \\
Catalan by DP                         & $O(n^2)$   & $O(n)$      & \coderef{c01-catalan} \\
\bottomrule
\end{tabular}
\end{center}

\begin{figure}[H]
\centering
\begin{tikzpicture}[diagram, node distance=6mm,
    q/.style={boxnode, fill=primaryblue!12, draw=primaryblue, text width=30mm},
    a/.style={boxnode, fill=successgreen!12, draw=successgreen, text width=34mm}]
  \node[q] (ask) {\textbf{What does the tree problem ask?}};
  \node[a, below left=9mm and 0mm of ask] (t1) {visit all nodes / compare shape \\ \footnotesize DFS or BFS (\secref{c01-what-is-tree})};
  \node[a, below=9mm of ask] (t2) {a value per node from root / leaves / count below \\ \footnotesize depth, height, size DFS (\secref{c01-anatomy})};
  \node[a, below right=9mm and 0mm of ask] (t3) {a shape promise? \\ \footnotesize full / complete / perfect / balanced (\secref{c01-kinds})};
  \node[a, right=3mm of t3, yshift=-14mm] (t4) {count trees / BSTs? \\ \footnotesize Catalan, Cayley (\secref{c01-counting})};
  \draw[flow] (ask) -- (t1); \draw[flow] (ask) -- (t2); \draw[flow] (ask) -- (t3); \draw[softarrow] (t3) -- (t4);
\end{tikzpicture}
\caption*{Choosing a Chapter-1 tool from how the problem is phrased. Harder phrasings
(path queries, subtree updates) point to later chapters.}
\end{figure}

\chaptersection{Common Mistakes Checklist}

\begin{itemize}[nosep, leftmargin=1.4em]
  \item \textbf{Missing the \nil{}/leaf base case.} Every tree recursion must say what a
        null (or a leaf) returns \emph{before} recursing. Fix: write the base case first.
  \item \textbf{Depth vs height, edges vs nodes.} They differ by one and run in opposite
        directions. Fix: pick the convention, set the base case to match, test a one-node tree.
  \item \textbf{Forgetting the parent guard.} A DFS on an adjacency list without
        \code{if (to == p) continue;} loops forever. Fix: always pass the parent.
  \item \textbf{Recursion on a skewed tree.} A path of $10^5$ nodes overflows the stack.
        Fix: iterative post-order with an explicit stack.
  \item \textbf{Full / complete / perfect mixed up.} They give different counting rules.
        Fix: sketch a two- and three-level example of each before asserting one.
  \item \textbf{Balanced checked only at the root.} The condition is per-node.
        Fix: return height bottom-up and a sentinel the moment any node breaks it.
  \item \textbf{Shapes vs labelled trees.} $C_n$ ignores labels; $n^{\,n-2}$ is unrooted
        and unbinary. Fix: ask whether identities matter and whether binary structure is required.
  \item \textbf{Overflow on counts.} Catalan and $n^{n-2}$ explode past tiny $n$.
        Fix: reduce modulo the given prime at every multiply.
\end{itemize}

\chaptersection{Quick Reference Card}

\begin{quickref}[Chapter 1 at a glance]
\textbf{Definitions.} Tree = connected + acyclic. Rooted: parent/child/leaf/subtree/ancestor.
$\depth$ down from root, $\height$ down to deepest leaf, $\size$ nodes in subtree.

\medskip
\textbf{Core identities.} $n$ nodes $\Rightarrow n-1$ edges. Unique path between any two nodes.
$\height(\nil){=}-1,\ \height(v){=}1{+}\max_c\height(c)$; $\size(v){=}1{+}\sum_c\size(c)$.
Perfect: $2^{h+1}{-}1$ nodes, $2^h$ leaves. Full binary: $L{=}I{+}1$.
Height: $\floor{\log_2 n}\le h\le n{-}1$.

\medskip
\textbf{Counting.} Binary shapes $C_n=\tfrac{1}{n+1}\binom{2n}{n}=1,1,2,5,14,42,\dots$;
labelled trees $n^{\,n-2}$; labelled rooted trees $n^{\,n-1}$.

\medskip
\textbf{Conventions.} Heights/depths in \emph{edges} unless a problem says nodes
(LeetCode \emph{Max Depth} counts nodes: subtract $1$ to convert).
\end{quickref}

\begin{lstlisting}[caption={Chapter 1 code kit: the two shapes you reuse everywhere}, label={lst:c01-kit}]
#include <bits/stdc++.h>
using namespace std;

// (1) General tree from an edge list: DFS with a parent guard, all three measures.
int n; vector<vector<int>> adj; vector<int> dep, hei, sz;
void measure(int v, int p, int d) {
    dep[v] = d; sz[v] = 1; hei[v] = -1;                 // edges convention
    for (int to : adj[v]) if (to != p) {
        measure(to, v, d + 1);
        sz[v]   += sz[to];
        hei[v]   = max(hei[v], hei[to]);
    }
    ++hei[v];
}

// (2) Binary tree height (edges) and node-count height (nodes) in two lines.
int heightEdges(TreeNode* r) { return !r ? -1 : 1 + max(heightEdges(r->left), heightEdges(r->right)); }
int heightNodes(TreeNode* r) { return !r ?  0 : 1 + max(heightNodes(r->left), heightNodes(r->right)); }
\end{lstlisting}

\chaptersection{Mixed Problem Set}

These blend the chapter's ideas and do not name the topic; decide the tool yourself.

\begin{mixedset}
  \practice{Merge Two Binary Trees}{LeetCode 617}
    {https://leetcode.com/problems/merge-two-binary-trees/}{LC Easy}
    {recurse in lockstep like Same Tree, but build a new node from the sum at each pair.}
  \practice{Path Sum}{LeetCode 112}
    {https://leetcode.com/problems/path-sum/}{LC Easy}
    {carry a running total down; a root-to-leaf path must hit the target exactly at a leaf.}
  \practice{Subtree of Another Tree}{LeetCode 572}
    {https://leetcode.com/problems/subtree-of-another-tree/}{LC Easy}
    {at each node try Same Tree; otherwise recurse into both children.}
  \practice{Lowest Common Ancestor of a BST}{LeetCode 235}
    {https://leetcode.com/problems/lowest-common-ancestor-of-a-binary-search-tree/}{LC Medium}
    {use the BST order: descend until the two targets split sides. See \chref{c11}.}
  \practice{Lowest Common Ancestor of a Binary Tree}{LeetCode 236}
    {https://leetcode.com/problems/lowest-common-ancestor-of-a-binary-tree/}{LC Medium}
    {Stretch: post-order returning which target was found; the node where both appear is the LCA (\chref{c11}).}
  \practicehh{Binary Tree Maximum Path Sum}{LeetCode 124}
    {https://leetcode.com/problems/binary-tree-maximum-path-sum/}{LC Hard}
    {Stretch: like the diameter -- a global best uses the best path THROUGH each node.}
    {DFS returns the best downward gain from a node ($\max(0, \text{child})$); update the answer with node + left + right.}
\end{mixedset}

\chaptersection{Related Topics}

\begin{itemize}[nosep, leftmargin=1.4em]
  \item \chref{c02} -- Binary trees in depth: the \code{TreeNode} model, construction, and more traversals.
  \item \chref{c03} -- Preorder, inorder, postorder and level order as reusable walks.
  \item \chref{c04} -- Storing trees: arrays, adjacency lists, parent/child pointers, Euler tour.
  \item \chref{c06} -- Self-balancing trees, where the height-bounds corollary becomes a design goal.
  \item \chref{c09} -- Diameter and centroid, built on the same bottom-up height idea.
  \item \chref{c10} -- Tree DP: the "combine children, return a value" pattern, formalised.
  \item \chref{c11} -- Lowest common ancestor, the path tool this chapter only hinted at.
  \item \chref{c17} -- Pr\"ufer codes: the full bijection behind Cayley's $n^{\,n-2}$.
\end{itemize}
